\section{シミュレーションにおける評価方法}
\label{sec:evaluation_method}

軟弱地面条件下における二足歩行の安定性を評価するため，シミュレーション結果から複数の力学的指標を抽出し，それらの挙動を解析した．評価指標としては，重心（Center of Mass, CoM）の挙動，ZMP（Zero Moment Point）の軌道，および関節トルクの時間変化に着目した．これらは，歩行安定性を異なる観点から評価可能な代表的指標である．

\subsection{重心（CoM）の挙動による評価}
\label{subsec:evaluation_com}

CoM の位置および高さの時間変化は，歩行中の全身運動の安定性を直接的に反映する指標である．特に軟弱地面においては，床沈み込みによりロボットの支持条件が変化するため，CoM の鉛直方向挙動が歩行安定性に大きな影響を及ぼす．

本研究では，歩行中の CoM の水平方向および鉛直方向の軌道を記録し，床沈み込みが発生した際の CoM の応答特性を解析することで，従来手法における安定性の変化を評価した．

\subsection{ZMP 軌道による評価}
\label{subsec:evaluation_zmp}

ZMP は，ロボットが床と接触している際の力学的安定性を表す指標であり，ZMP が支持多角形内に存在するかどうかは，静的および準動的な安定性評価において重要な基準となる．

本研究では，歩行中に算出される ZMP の時間軌道を解析し，床沈み込みが生じた際に ZMP が支持脚の足裏領域内でどのように変化するかを評価した．

\subsection{関節トルク挙動による評価}
\label{subsec:evaluation_torque}

関節トルクは，歩行制御系が要求される力学的負荷を直接的に表す指標である．軟弱地面では，床沈み込みにより想定外の姿勢変化が生じるため，特定の関節に過大なトルクが要求される可能性がある．

本研究では，特に支持脚のロール方向関節トルクに着目し，両脚支持期から片脚支持期への遷移時におけるトルクの時間変化を解析した．過大なトルクの発生や急峻なピークは，転倒や制御破綻の要因となるため，歩行安定性を評価する上で重要な指標として扱った．

以上の指標を総合的に解析することで，軟弱地面条件下における従来手法の歩行安定性を評価するとともに，不安定化要因の特定を行った．次章では，これらの評価結果に基づき，軟弱地面歩行における不安定挙動の特徴について詳しく述べる．
